\stagechapter{R}{Reporting}

\noindent\textbf{What this stage covers.} What the thesis has to report so that a reader can follow what was planned, done and found: the search and selection numbers, the capture records, the codebook and its versions, reliability, the LLM configuration, the harm tree and the limits. No single guideline fits an app-and-policy mapping study, so the corpus offers a set of partial ones that can be combined, and the thesis should say which items it adapted rather than claim compliance. Draft plan: results are descriptive and reported per domain, denominators are explicit, absence means ``no mention identified in the captured material'', and the limitations section states what the study cannot establish.

\section{Reporting guidelines in the corpus}
\begin{table}[t]
\centering\small
\caption{Guidelines and the parts of a thesis they support.}
\label{tab:R-guidelines}
\begin{tabularx}{\textwidth}{@{}P{2.7cm}P{3.3cm}Y@{}}
\toprule
\textbf{Guideline} & \textbf{Scope} & \textbf{Items to adapt}\\
\midrule
\pa{page2021} PRISMA 2020 & 27 items for systematic reviews; flow diagram & eligibility, sources and dates, search, who screened and how independently, numbers with reasons for exclusion, data items, availability of data, code and materials\\
\pa{rethlefsen2021} PRISMA-S & 16 items for the search & platform, strategies as run, limits with justification, dates, record counts, deduplication; web and registry searches\\
\pa{tricco2018} PRISMA-ScR & 20 items plus 2 optional for scoping reviews & rationale for scoping, ``sources of evidence'' versus ``information sources'', data charting, no risk-of-bias scoring, funding of included sources\\
\pa{arksey2005} & conduct framework for scoping studies & question, comprehensive search, post hoc criteria logged, charting form, no synthesis, template per group\\
\pa{vonelm2007} STROBE & 22 items for cohort, case-control, cross-sectional studies & setting and dates, eligibility, variables and measurement, bias, study size, missing data, flow of numbers, limitations, generalisability\\
\pa{obrien2014} SRQR & 21 items for qualitative research & paradigm and approach, researcher characteristics, sampling, unit of analysis, instruments, trustworthiness, quotes as evidence\\
\bottomrule
\end{tabularx}
\end{table}
All six say that they guide reporting, not conduct, and none is a quality score: PRISMA 2020 ``should not be used to assess the conduct or methodological quality of systematic reviews'' \pc{page2021}{5}{PRISMA 2020 should not be used to assess the conduct or methodological quality of systematic reviews}; STROBE ``was not developed as a tool for assessing the quality of published observational research'' \pc{vonelm2007}{4}{We emphasise that the STROBE Statement was not developed as a tool for assessing the quality of published observational research} and warns against ``STROBEing'' a manuscript; PRISMA-S wants strategies ``copied and pasted exactly as run'' \pc{rethlefsen2021}{4}{Include the search strategies for each database and information source, copied and pasted exactly as run}; PRISMA-ScR treats the ``charting'' form and its independence as reportable \pc{tricco2018}{24}{whether data charting was done independently, in duplicate}.

\section{Reporting the analysis}
\begin{itemize}
\item \emph{Qualitative coding.} \pa{obrien2014} asks for the unit of analysis \pc{obrien2014}{29}{Authors should specify the unit of analysis}, how discrepant cases were handled \pc{obrien2014}{30}{the authors must describe how these discrepancies were handled during the analysis}, the software used (``simply stating that software was used is insufficient'') and an audit trail \pc{obrien2014}{33}{An audit trail involves careful documentation of all decisions made throughout the study}. \pa{nowell2017} gives a per-phase trustworthiness checklist (audit trail, versioned codebook, reflexive journal, minutes of coder discussions). \pa{braun2021} offers a 20-question evaluation tool and warns against ten misuses; \pa{braun2006} lists reporting items (state assumptions, describe what was done, avoid ``themes emerged'').
\item \emph{Reliability.} \pa{lombard2002} lists ten items (pilot size, reliability sample, independence, a second index beside percent agreement, per-code values, training, resolution rule); \pa{hallgren2012} adds design, metric, marginals, the variant of the statistic and the effect on later analyses; \pa{mcdonald2019} asks for thresholds per stage, even pre-registered; \pa{oconnor2020} for per-code distributions. See the checklist in chapter~I.
\item \emph{LLM use.} \pa{tornberg2024}, \pa{baumann2025} and \pa{pangakis2023}: model version, prompt, parameters, all configurations tried, validation set size and metrics per code; see chapter~L.
\item \emph{Harm trees.} \pa{knabenhans2026} on presenting trees with origins and the statement that completeness is not claimed.
\item \emph{Ethics and data.} \pa{aoir2019}: stage-based ethical reflection, data minimisation, storage and retention, terms of platforms for automated collection.
\end{itemize}

\section{Reporting the limits}
Several corpus papers model careful limitation wording: \pa{magoulas2026} describes store disclosures as presentations, not verified practice; \pa{tangari2021} and \pa{arkalakis2024} show what behaviour analysis adds; \pa{malki2024} lists a single platform, single-inspector usability inspection and no legal specialists; \pa{doan2025} flags a single coder as a reliability gap; \pa{hutton2018} separates popularity from privacy quality. \pa{vonelm2007} asks for the direction and magnitude of possible bias and for generalisability; \pa{tricco2018} and \pa{arksey2005} for a statement that quality was not appraised.

\begin{table}[t]
\centering\small
\caption{Where each element is reported in a thesis of this design (proposal).}
\label{tab:R-map}
\begin{tabularx}{\textwidth}{@{}P{3.1cm}YP{3.3cm}@{}}
\toprule
\textbf{Element} & \textbf{Content} & \textbf{Guidance}\\
\midrule
Search log (stores) & storefront, date, device, account, keywords, results examined, filters, de-duplication & \pa{rethlefsen2021}, \pa{page2021}\\
Flow of apps & identified, screened, excluded with reasons, included; pilot apps flagged & \pa{page2021}, \pa{vonelm2007}\\
Capture records & URL, date, language, hash, app version; shared policies linked & \pa{shipp2020}, \pa{light2018}\\
Codebook & definitions, examples, counter-examples, versions, change log, release date & \pa{macqueen1998}, \pa{nowell2017}\\
Reliability & subset, coders, statistic, per-code values with marginals, thresholds, resolution & \pa{lombard2002}, \pa{hallgren2012}\\
LLM configuration & model, version, prompts, parameters, all configurations tried, validation metrics & \pa{baumann2025}, \pa{tornberg2024}\\
Harm trees & structure, sources, origin tags, applicability record & \pa{knabenhans2026}\\
Results & denominators, absence wording, counts per domain and surface & \pa{vonelm2007}, \pa{arksey2005}\\
Limits & bias direction, generalisability, what cannot be claimed & \pa{vonelm2007}, \pa{magoulas2026}\\
Ethics, data & terms checked, retention, personal data excluded & \pa{aoir2019}\\
\bottomrule
\end{tabularx}
\end{table}

\begin{thesisbox}[Implications for this thesis (stage R)]
Do not write that the thesis ``follows PRISMA'' or ``adheres to STROBE''. State instead that the review and selection parts adapt PRISMA-ScR, PRISMA-S and the PRISMA flow diagram, that the empirical description adapts STROBE items for a cross-sectional study of apps, and that the qualitative coding follows SRQR items, listing the items not used and why. Provide an availability statement for the codebook, prompts, extraction tables and code (PRISMA 2020 item 27), and keep the Decision Log as the audit trail for deviations (PRISMA 2020 item 24c).
\end{thesisbox}
